\subsection{Messverfahren}
\xfigure{htbp}{scale=1}{Versuchsaufbau.png}{Versuchsaufbau: Prismenspektrometer \autocite{Skript_1}}{Prismenspektrometer \autocite{Skript_1}}

Durch das Messen der Ablenkwinkel \(\delta\) der Emissionslinien einer Helium(\ce{He})-Lampe wird die Kurve \(\delta\ (\lambda)\) des Prismas bestimmt, also fast \(n(\lambda)\), da \(\delta = \delta (n(\lambda))\).\\ Dafür nutzen wir das Prismenspektrometer, bestehend aus Kollimator-Rohr, welches das Gaslampenlicht parallelisiert; einem Prismentisch; Fernrohr und Nonius, mit dem man den Ablenkwinkel nach Ausrichten des Fernrohres auf eine Spektrallinie ablesen kann.

Nun werden die Ablenkwinkel der Emissionslinien einer Quecksilber-(\ce{Hg}) und einer Wasserstofflampe (\ce{H}) gemessen, und mit erstellter Eichkurve \(\delta\ (\lambda)\) werden die Wellenlängen der entsprechenden Linien bestimmt und mit Literaturwerten verglichen.
\subsection{Messprotokoll}
\includegraphics[page=1,width=\textwidth]{V33_12.pdf}
\includegraphics[page=2,width=\textwidth]{V33_12.pdf}
\newpage
\includegraphics[page=1,width=\textwidth]{V33_3.pdf}